% Generated draft: requires mathematical and source review.
\documentclass[11pt]{article}
\usepackage[margin=25mm]{geometry}
\usepackage{fontspec}
\setmainfont{Latin Modern Roman}
\usepackage{amsmath,amssymb,mathtools,unicode-math}
\setmathfont{Latin Modern Math}
\usepackage{enumitem,needspace,xurl,hyperref}
\setlength{\parindent}{0pt}
\setlength{\parskip}{5pt}
\setlength{\emergencystretch}{3em}
\newcommand{\N}{\mathbb{N}}
\newcommand{\Z}{\mathbb{Z}}
\newcommand{\Q}{\mathbb{Q}}
\newcommand{\R}{\mathbb{R}}
\newcommand{\C}{\mathbb{C}}
\newcommand{\F}{\mathbb{F}}
\newcommand{\A}{\mathbb{A}}
\newcommand{\PP}{\mathbb{P}}
\newcommand{\E}{\mathbb{E}}
\newcommand{\eps}{\varepsilon}
\newcommand{\dd}{\,\mathrm{d}}
\begin{document}
\section*{UnsolvedMath ID: 800940}
\textbf{Source problem code:} catalogue entry
\subsection{Definitions and mathematical statement}
Let \((M,g)\) be a closed Riemannian manifold, meaning compact and without boundary, of dimension \(n>2\), with sectional curvature \(K_g(\sigma)<0\) for every tangent two-plane \(\sigma\). Write
\[
SM=\{v\in TM:g(v,v)=1\}
\]
for its unit tangent bundle and \(\varphi_t:SM\to SM\) for its geodesic flow. Its Anosov splitting is
\[
T(SM)=E^s\oplus\R X\oplus E^u,
\]
where \(X\) generates the flow, the subbundles are invariant under \(D\varphi_t\), and there exist constants \(C\geq1\) and \(\lambda>0\) such that, for every \(t\geq0\),
\[
\|D\varphi_t\xi\|\leq Ce^{-\lambda t}\|\xi\|
\quad(\xi\in E^s),\qquad
\|D\varphi_{-t}\eta\|\leq Ce^{-\lambda t}\|\eta\|
\quad(\eta\in E^u).
\]
Here the norms may be taken from the Sasaki metric on \(SM\). Saying that \(E^s\) and \(E^u\) are \(C^1\) means that they are \(C^1\) subbundles over all of \(SM\), equivalently that their associated sections of the Grassmann bundle are \(C^1\).

\textbf{Existence question (Burns--Matveev, Question 20).} Does there exist such a manifold \((M,g)\) for which both \(E^s\) and \(E^u\) are \(C^1\), but \(g\) is not \(1/4\)-pinched?

\textbf{Unresolved pinching qualification.} The supplied record does not define \(1/4\)-pinching, so the last condition cannot be given a unique numerical interpretation from this record alone. For a global curvature interpretation, define
\[
\kappa_{\min}=\min_{\sigma}(-K_g(\sigma)),\qquad
\kappa_{\max}=\max_{\sigma}(-K_g(\sigma)),\qquad
r(g)=\frac{\kappa_{\min}}{\kappa_{\max}}.
\]
These extrema range over all tangent two-planes, and \(0<\kappa_{\min}\leq\kappa_{\max}\). If \(1/4\)-pinching means \(r(g)\geq1/4\), the requested failure is \(r(g)<1/4\); if it means \(r(g)>1/4\), the requested failure is \(r(g)\leq1/4\). A pointwise pinching interpretation would instead compare curvature extrema separately at each base point. The supplied evidence does not select among these interpretations.

\textbf{Supplied context.} The record states that the splitting is at least \(C^1\) in dimension two or under \(1/4\)-pinching, and that in higher dimensions it is always H\"older continuous but typically not \(C^1\). It refers to the introduction of \([\mathrm{Hasselblatt1994}]\) for regularity results. These are supplied contextual assertions; the pinching convention and the meaning of ``typically'' are not specified here.
\subsection{Short English statement}
Burns and Matveev ask whether a closed negatively curved Riemannian manifold of dimension greater than two can have both stable and unstable geodesic-flow distributions of class \(C^1\) while failing \(1/4\)-pinching. The supplied record leaves the pinching convention, including its equality boundary, unresolved. Its registry status is supplied as open; that status has not been independently verified.
\subsection{Sources}
\begin{itemize}
\item Local UnsolvedMath registry, permanent ID \(800940\), exact source problem code \texttt{catalogue entry}, titled \emph{Open Problems and Questions about Geodesics --- Question 20}. Supplied attribution: Burns, Matveev; proposed year: 2019; supplied status: open. The record belongs to the AMR Open Problem Lists collection.
\item Burns and Matveev, \emph{Open Problems and Questions about Geodesics}, Question 20; the supplied source list labels the work 2019. Supplied source URL: \url{https://arxiv.org/abs/1308.5417}. The provenance records access on 2026-10-16 and extraction by \texttt{arxiv\_latex}; these are supplied provenance details, not a claim of independent access or verification.
\item The statement cites \([\mathrm{Hasselblatt1994}]\), specifically its introduction, for a list of regularity results. No full bibliographic entry or link for this reference is supplied.
\item Supplied upstream revision: \texttt{d371ee917ac9647fbb06cd6fc529deec6f3dda82}. The registry record was created and updated on 2026-10-16. These dates do not establish a current literature assessment.
\end{itemize}
\end{document}
